\documentclass[../main.tex]{subfiles}
\begin{document}
本章主要参考\cite{brezis,hunter}。

% Reference: https://ocw.mit.edu/courses/18-125-measure-and-integration-fall-2003/pages/lecture-notes/
\section{测度与可测函数}
\subsection{测度空间与Lebesgue测度}\index{测度空间与Lebesgue测度|hyperpage}
测度空间记为 $(X,\mathcal A,\mu)$，其中 $\mathcal A$ 为 $\sigma$-代数，$\mu(\varnothing)=0$，且对两两不交的 $E_n\in\mathcal A$，
\[
\mu\!\left(\bigcup_{n=1}^\infty E_n\right)=\sum_{n=1}^\infty\mu(E_n).
\]
$\mathbb R^d$ 上 Lebesgue 测度满足矩形体积公式 $m(\prod_{j=1}^d[a_j,b_j])=\prod_j(b_j-a_j)$。性质在零测集之外成立称为几乎处处成立，记 a.e.。
\subsection{测度的连续性与可测性}\index{测度的连续性与可测性|hyperpage}
\[
E_n\uparrow E\Longrightarrow\mu(E_n)\to\mu(E),\qquad
E_n\downarrow E,\ \mu(E_1)<\infty\Longrightarrow\mu(E_n)\to\mu(E).
\]
实值函数 $f$ 可测当且仅当 $\{x:f(x)>a\}\in\mathcal A$ 对每个 $a\in\mathbb R$ 成立；可测函数的逐点极限仍可测。第二个测度连续性公式不可省略有限测度条件。
\section{Lebesgue积分}
\subsection{定义与基本不等式}\index{定义与基本不等式|hyperpage}
对非负简单函数 $s=\sum_{k=1}^na_k\bm1_{E_k}$，$E_k$ 两两不交，定义 $\int s\,\mathrm d\mu=\sum_ka_k\mu(E_k)$（约定 $0\cdot\infty=0$）。对非负可测 $f$，
\[
\int f\,\mathrm d\mu=\sup_{0\le s\le f,\ s\text{ 为简单函数}}\int s\,\mathrm d\mu.
\]
一般实值函数分解为 $f=f^+-f^-$；若 $\int|f|\,\mathrm d\mu<\infty$，称 $f\in L^1$，且
\[
\int f=\int f^+-\int f^-,\qquad\left|\int f\right|\le\int|f|.
\]
\subsection{单调收敛、Fatou引理与控制收敛}\index{单调收敛、Fatou引理与控制收敛|hyperpage}
对非负可测函数，单调收敛与 Fatou 引理为
\[
0\le f_n\uparrow f\Longrightarrow\int f_n\to\int f,\qquad
\int\liminf_nf_n\le\liminf_n\int f_n.
\]
若 $f_n\to f$ a.e. 且存在 $g\in L^1$ 使 $|f_n|\le g$ a.e.，则控制收敛定理给出
\[
\int|f_n-f|\to0,\qquad\int f_n\to\int f.
\]
逐点收敛本身不足以交换极限与积分；需检查单调性、控制函数或其他充分条件。
\subsection{Tonelli与Fubini定理}\index{Tonelli与Fubini定理|hyperpage}
在两个 $\sigma$-有限测度空间的乘积上，非负可测函数满足 Tonelli 定理，允许积分值为 $+\infty$：
\[
\int_{X\times Y}f\,\mathrm d(\mu\times\nu)
=\int_X\!\left(\int_Yf\,\mathrm d\nu\right)\mathrm d\mu
=\int_Y\!\left(\int_Xf\,\mathrm d\mu\right)\mathrm d\nu.
\]
若 $f$ 可测且 $\int_{X\times Y}|f|<\infty$，则 Fubini 定理保证两种累次积分几乎处处有定义、可积且等于乘积积分。
\section{函数空间与收敛}
\subsection{Lp范数与基本不等式}\index{Lp范数与基本不等式|hyperpage}
对 $1\le p<\infty$，$\|f\|_p=(\int|f|^p\,\mathrm d\mu)^{1/p}$；$\|f\|_\infty=\operatorname*{ess\,sup}|f|$。$L^p$ 中将 a.e. 相等的函数视为同一元素。
\[
\int|fg|\le\|f\|_p\|g\|_q\quad(1/p+1/q=1),\qquad
\|f+g\|_p\le\|f\|_p+\|g\|_p.
\]
这里 $1\le p,q\le\infty$，约定 $1/\infty=0$。若 $0<\mu(X)<\infty$、$1\le p\le q\le\infty$，则
\[
\|f\|_p\le\mu(X)^{1/p-1/q}\|f\|_q.
\]
\subsection{几种收敛方式}\index{几种收敛方式|hyperpage}
\begin{align*}
f_n\to f\text{ 依测度： }&\quad\forall\varepsilon>0,\ \mu(|f_n-f|>\varepsilon)\to0,\\
f_n\to f\text{ 于 }L^p\text{：}&\quad\|f_n-f\|_p\to0,\\
f_n\to f\text{ 一致： }&\quad\sup_x|f_n(x)-f(x)|\to0.
\end{align*}
对 $p<\infty$，$L^p$ 收敛蕴含依测度收敛，因为
\[
\mu(|f_n-f|>\varepsilon)\le\varepsilon^{-p}\|f_n-f\|_p^p.
\]
依测度收敛可选取 a.e. 收敛子列；在有限测度空间上，a.e. 收敛蕴含依测度收敛。
\section{绝对连续与微分}
\subsection{绝对连续函数与积分基本定理}\index{绝对连续函数与积分基本定理|hyperpage}
$f\in AC[a,b]$ 当且仅当存在 $g\in L^1[a,b]$ 使
\[
f(x)=f(a)+\int_a^xg(t)\,\mathrm dt.
\]
此时 $f'=g$ a.e.，$f(b)-f(a)=\int_a^bf'(x)\,\mathrm dx$。Lipschitz 函数绝对连续，绝对连续函数有界变差；仅连续且 a.e. 可导不足以推出此积分公式。
\subsection{Radon--Nikodym导数}\index{Radon-Nikodym导数|hyperpage}
若 $\mu,\nu$ 是同一可测空间上的 $\sigma$-有限正测度，且 $\nu\ll\mu$（$\mu(E)=0\Rightarrow\nu(E)=0$），则存在非负可测 $h$，使
\[
\nu(E)=\int_Eh\,\mathrm d\mu,\qquad h=\frac{\mathrm d\nu}{\mathrm d\mu}.
\]
该导数在 $\mu$-a.e. 意义下唯一；对可积 $f$，$\int f\,\mathrm d\nu=\int fh\,\mathrm d\mu$。
% Further references: MIT 18.125 notes; Oxford Function Spaces and Distribution Theory.
\section{测度构造与分解补充}
\subsection{外测度、Carathéodory条件与完备化}\index{外测度、Carathéodory条件与完备化|hyperpage}
由可数覆盖定义的外测度满足
\[
\mu^*(E)=\inf\left\{\sum_j\mu_0(A_j):E\subseteq\bigcup_jA_j\right\},\qquad
\mu^*\!\left(\bigcup_jE_j\right)\le\sum_j\mu^*(E_j).
\]
集合 $A$ 为 Carathéodory 可测，当且仅当对每个 $E\subseteq X$ 有
\[
\mu^*(E)=\mu^*(E\cap A)+\mu^*(E\setminus A).
\]
可测集构成 $\sigma$-代数，限制测度完备；$\sigma$-有限预测度在所生成的 $\sigma$-代数上延拓唯一。原测度的完备化加入全部零测集的子集，满足 $E=A\cup N$、$N\subseteq Z$、$\mu(Z)=0$ 时 $\bar\mu(E)=\mu(A)$。
\subsection{正则性、推前与变量替换}\index{正则性、推前与变量替换|hyperpage}
Lebesgue 可测集 $E\subseteq\mathbb R^d$ 满足外正则性；若 $m(E)<\infty$，还有内正则性：
\[
m(E)=\inf_{U\supseteq E,\ U\text{ 开}}m(U)
=\sup_{K\subseteq E,\ K\text{ 紧}}m(K).
\]
可测映射 $T:X\to Y$ 的推前测度为 $T_\#\mu(B)=\mu(T^{-1}B)$，且对非负或可积 $g$，
\[
\int_Yg\,\mathrm d(T_\#\mu)=\int_Xg\circ T\,\mathrm d\mu.
\]
若 $T:U\to V$ 为开集间 $C^1$ 微分同胚，则
\[
\int_Vg(y)\,\mathrm dy=\int_Ug(T(x))|\det DT(x)|\,\mathrm dx.
\]
\subsection{符号测度、全变差与Lebesgue分解}\index{符号测度、全变差与Lebesgue分解|hyperpage}
有限实符号测度有唯一 Jordan 分解 $\nu=\nu^+-\nu^-$，$\nu^+\perp\nu^-$。全变差为
\[
|\nu|(E)=\sup_{E=\bigsqcup_jE_j}\sum_j|\nu(E_j)|=\nu^+(E)+\nu^-(E),\qquad
\left|\int f\,\mathrm d\nu\right|\le\int|f|\,\mathrm d|\nu|.
\]
对 $\sigma$-有限正测度 $\mu$，有限符号或复测度可唯一分解为
\[
\nu=h\mu+\nu_s,\qquad h\in L^1(\mu),\quad \nu_s\perp\mu.
\]
若 $\nu=h\mu$，则 $|\nu|=|h|\mu$。对链式绝对连续的 $\sigma$-有限正测度，
\[
\frac{\mathrm d\lambda}{\mathrm d\mu}
=\frac{\mathrm d\lambda}{\mathrm d\nu}\frac{\mathrm d\nu}{\mathrm d\mu}\quad\mu\text{-a.e.}
\]
\section{积分与收敛的进一步判别}
\subsection{分布函数、层蛋糕与重排}\index{分布函数、层蛋糕与重排|hyperpage}
对可测 $f$ 定义 $d_f(t)=\mu(|f|>t)$。Tonelli 定理给出（$p>0$）
\[
\int|f|^p\,\mathrm d\mu=p\int_0^\infty t^{p-1}d_f(t)\,\mathrm dt.
\]
递减重排 $f^*(s)=\inf\{t\ge0:d_f(t)\le s\}$ 与 $|f|$ 等分布，故 $\|f\|_p^p=\int_0^\infty f^*(s)^p\,\mathrm ds$。Hardy--Littlewood 不等式为
\[
\int|fg|\,\mathrm d\mu\le\int_0^\infty f^*(s)g^*(s)\,\mathrm ds.
\]
弱 $L^p$ 量定义为 $\|f\|_{p,\infty}=\sup_{t>0}t\,d_f(t)^{1/p}$，它通常只是等价于某些范数的拟范数量。
\subsection{Egorov、Lusin与一致可积性}\index{Egorov、Lusin与一致可积性|hyperpage}
在有限测度空间上，$f_n\to f$ a.e. 蕴含：任意 $\varepsilon>0$，存在 $\mu(E)<\varepsilon$，使在 $X\setminus E$ 上一致收敛。Lebesgue 可测有限值函数在有限测度可测集上可去掉任意小测度集合而变为连续函数（Lusin 定理）。

有限测度空间中的 $\mathcal F\subset L^1$ 一致可积可写为
\[
\lim_{K\to\infty}\sup_{f\in\mathcal F}\int_{|f|>K}|f|\,\mathrm d\mu=0.
\]
Vitali 判别：若 $f_n\to f$ 依测度，则 $\|f_n-f\|_1\to0$ 当且仅当 $\{f_n\}$ 一致可积。充分条件为 $\sup_n\|f_n\|_p<\infty$（某个 $p>1$）；更一般地，若递增凸函数 $\Phi$ 满足 $\Phi(t)/t\to\infty$ 且 $\sup_n\int\Phi(|f_n|)<\infty$，也可保证一致可积。
\subsection{Scheffé引理与积分号下微分}\index{Scheffé引理与积分号下微分|hyperpage}
若 $f_n,f\ge0$、$f_n\to f$ a.e. 且 $\int f_n\to\int f<\infty$，则
\[
\|f_n-f\|_1=\int f_n+\int f-2\int\min(f_n,f)\longrightarrow0.
\]
设 $f(x,t)$ 对 $t$ 在开区间可微，$f(\cdot,t_0)\in L^1$，且每个紧参数区间内 $|\partial_tf(x,t)|\le g(x)$、$g\in L^1$，则
\[
\frac{\mathrm d}{\mathrm dt}\int f(x,t)\,\mathrm d\mu(x)
=\int\partial_tf(x,t)\,\mathrm d\mu(x).
\]
\subsection{密度、平移与磨光}\index{密度、平移与磨光|hyperpage}
在 $\mathbb R^d$ 上，$1\le p<\infty$ 时 $C_c^\infty$ 在 $L^p$ 中稠密，且
\[
\|f(\cdot-h)-f\|_p\to0\quad(h\to0).
\]
取非负 $\rho\in C_c^\infty$、$\int\rho=1$，$\rho_\varepsilon(x)=\varepsilon^{-d}\rho(x/\varepsilon)$，则
\[
f*\rho_\varepsilon\to f\text{ 于 }L^p,\qquad
\partial^\alpha(f*\rho_\varepsilon)=f*(\partial^\alpha\rho_\varepsilon).
\]
Young 不等式为 $\|f*g\|_r\le\|f\|_p\|g\|_q$，其中 $p,q,r\ge1$ 且 $1+1/r=1/p+1/q$。一般 $L^\infty$ 函数不具有范数平移连续性。
\section{微分、变差与弱导数补充}
\subsection{极大函数与Lebesgue微分定理}\index{极大函数与Lebesgue微分定理|hyperpage}
对 $f\in L^1_{\rm loc}(\mathbb R^d)$ 定义中心极大函数
\[
Mf(x)=\sup_{r>0}\frac1{|B(x,r)|}\int_{B(x,r)}|f(y)|\,\mathrm dy.
\]
有弱 $(1,1)$ 估计 $|\{Mf>\lambda\}|\le C_d\|f\|_1/\lambda$；当 $1<p\le\infty$ 时 $\|Mf\|_p\le C_{d,p}\|f\|_p$。几乎每个 $x$ 为 Lebesgue 点：
\[
\lim_{r\downarrow0}\frac1{|B(x,r)|}\int_{B(x,r)}|f(y)-f(x)|\,\mathrm dy=0.
\]
因此可用局部平均恢复函数的几乎处处代表。
\subsection{有界变差、Stieltjes测度与分部积分}\index{有界变差、Stieltjes测度与分部积分|hyperpage}
\[
\operatorname{Var}_{[a,b]}f=\sup_{a=x_0<\cdots<x_n=b}\sum_i|f(x_i)-f(x_{i-1})|.
\]
实函数有界变差当且仅当可表示成两个递增函数之差；此时 a.e. 可导且 $\int_a^b|f'|\le\operatorname{Var}f$。右连续递增 $F$ 诱导测度 $\mu_F((s,t])=F(t)-F(s)$，其原子质量为 $F(x)-F(x-)$。若 $f,g\in AC[a,b]$，
\[
\int_a^bf g'=[fg]_a^b-\int_a^bf'g.
\]
对 $f\in BV(a,b)$，分布导数为有限符号测度，且 $|Df|(a,b)$ 为区间内总变差；$Df=f'\,\mathrm dx+D^sf$。绝对连续恰好要求奇异部分为零（采用相应连续代表）。
\subsection{弱导数与Sobolev基本公式}\index{弱导数与Sobolev基本公式|hyperpage}
在开集 $\Omega\subset\mathbb R^d$ 上，弱导数由
\[
\int_\Omega u\,\partial^\alpha\varphi
=(-1)^{|\alpha|}\int_\Omega D^\alpha u\,\varphi\quad(\varphi\in C_c^\infty(\Omega))
\]
定义。令 $\|u\|_{W^{k,p}}=(\sum_{|\alpha|\le k}\|D^\alpha u\|_p^p)^{1/p}$（$p<\infty$）。乘积与链式法则在各项可积且函数足够有界等条件下成立；例如 $u\in W^{1,p}$、$F\in C^1$ 且 $F'$ 有界时，$D(F(u))=F'(u)Du$ 局部成立。

对 $u\in C_c^\infty(\mathbb R^d)$、$1\le p<d$，
\[
\|u\|_{p^*}\le C\|\nabla u\|_p,\qquad p^*=\frac{dp}{d-p}.
\]
有界 Lipschitz 域上 $u\in W_0^{1,p}(\Omega)$ 满足 Poincaré 不等式 $\|u\|_p\le C_\Omega\|\nabla u\|_p$。当 $p>d$ 时局部有 Hölder 连续代表，指数 $1-d/p$；有界 Lipschitz 域上 $W^{1,p}\hookrightarrow L^q$ 对 $q<p^*$ 为紧嵌入（$p<d$），临界指数处一般不紧。
\subsection{Hausdorff测度与维数}\index{Hausdorff测度与维数|hyperpage}
对 $s\ge0$ 定义未归一化 Hausdorff 测度
\[
\mathcal H^s(E)=\lim_{\delta\downarrow0}\inf\left\{\sum_i(\operatorname{diam}U_i)^s:
E\subseteq\bigcup_iU_i,\ \operatorname{diam}U_i\le\delta\right\}.
\]
\[
\dim_HE=\inf\{s:\mathcal H^s(E)=0\},\qquad
\mathcal H^s(T(E))\le L^s\mathcal H^s(E)
\]
当 $T$ 为 $L$-Lipschitz 时成立，故 Lipschitz 映射不增 Hausdorff 维数，双 Lipschitz 映射保持维数。中间三分 Cantor 集的维数为 $\log2/\log3$。
% 联动扩充：实分析。References:
% https://math.ucdavis.edu/~hunter/pdes/pde_notes.pdf
% https://arxiv.org/abs/1104.4345
% https://sites.math.rutgers.edu/~brezis/PUBlications/233.pdf
\section{覆盖工具与精细微分}\label{sec:ra-covering}
\subsection{五倍覆盖引理与Vitali覆盖}\index{五倍覆盖引理与Vitali覆盖|hyperpage}
欧氏空间中半径一致有界的球族可选取两两不交的至多可数子族 $B_j$，使
\[
\bigcup_{B\in\mathcal B}B\subseteq\bigcup_j5B_j.
\]
若 $E$ 可测且测度有限，覆盖族在每个 $x\in E$ 提供任意小的、以 $x$ 为中心的球，则 Vitali 定理可选取不交球，使 $|E\setminus\bigcup_jB_j|=0$。后一结论使用细覆盖假设，不能由任意一次球覆盖直接推出。
\subsection{Besicovitch覆盖与有界重叠}\index{Besicovitch覆盖与有界重叠|hyperpage}
对有界 $E\subset\mathbb R^d$ 的中心细球覆盖，可选取至多可数球族覆盖 $E$，并满足
\[
\sum_j\bm1_{B_j}(x)\le N_d,
\qquad \sum_j\int_{B_j}|f|\le N_d\int_{\mathbb R^d}|f|.
\]
$N_d$ 仅依赖维数；也可按有限个不交球族组织覆盖。该几何性质不能无条件推广到任意度量空间，常用于测度微分及密度估计。
\subsection{密度点与本质边界}\index{密度点与本质边界|hyperpage}
对 Lebesgue 可测集 $E$，几乎每个点满足
\[
\lim_{r\downarrow0}\frac{|E\cap B(x,r)|}{|B(x,r)|}=\bm1_E(x).
\]
记密度为 $1,0$ 的点集为 $E^{(1)},E^{(0)}$，本质边界为 $\partial^eE=\mathbb R^d\setminus(E^{(1)}\cup E^{(0)})$，于是 $|\partial^eE|=0$。它不同于拓扑边界；后者可能有正体积。
\subsection{Rademacher可微性与Lipschitz微分}\index{Rademacher可微性与Lipschitz微分|hyperpage}
Lipschitz 映射 $F:\mathbb R^d\to\mathbb R^m$ 几乎处处具有全微分，且
\[
F(x+h)=F(x)+DF(x)h+o(|h|),\qquad\|DF(x)\|_{\rm op}\le\operatorname{Lip}(F).
\]
标量情形 $F\in W^{1,\infty}_{\rm loc}$，经典导数与弱导数 a.e. 一致。仅有各坐标偏导数存在，不足以推出全微分存在；面积与余面积公式见\hyperref[chap:measure-theory]{测度论}。
\section{Lp估计与定量收敛}
\subsection{积分型Minkowski与范数插值}\index{积分型Minkowski与范数插值|hyperpage}
对联合可测 $f(x,y)$、$1\le p\le\infty$，在右端有限的条件下，
\[
\left\|\int f(\cdot,y)\,\mathrm d\nu(y)\right\|_{L^p(\mu)}
\le\int\|f(\cdot,y)\|_{L^p(\mu)}\,\mathrm d\nu(y).
\]
在通常的 $\sigma$-有限乘积空间上可使用此式。若 $1/r=\theta/p+(1-\theta)/q$、$0\le\theta\le1$，则 $\|u\|_r\le\|u\|_p^\theta\|u\|_q^{1-\theta}$，其中 $1\le p,q\le\infty$。
\subsection{Clarkson不等式与强收敛判据}\index{Clarkson不等式与强收敛判据|hyperpage}
对 $2\le p<\infty$，
\[
\left\|\frac{f+g}2\right\|_p^p+\left\|\frac{f-g}2\right\|_p^p
\le\frac{\|f\|_p^p+\|g\|_p^p}2.
\]
对 $1<p\le2$、$q=p/(p-1)$，左端的两个幂改为 $q$，右端改为 $[(\|f\|_p^p+\|g\|_p^p)/2]^{q/p}$。这些估计给出一致凸性及 Radon--Riesz 性质：$f_n\rightharpoonup f$ 且 $\|f_n\|_p\to\|f\|_p$ 蕴含 $\|f_n-f\|_p\to0$（$1<p<\infty$）。
\subsection{Brezis--Lieb积分分解}\index{Brezis-Lieb积分分解|hyperpage}
若 $0<p<\infty$，$f_n\to f$ a.e. 且 $\sup_n\int|f_n|^p<\infty$，则 $f\in L^p$，并有
\[
\int|f_n|^p-\int|f_n-f|^p\longrightarrow\int|f|^p.
\]
此式允许把极限主体与剩余部分的 $p$ 次质量分离；不要求 $f_n\to f$ 于 $L^p$。$p<1$ 时相关量为拟范数，不能据此调用 Banach 空间结论。
\subsection{Lp型Vitali判据与无穷远紧性}\index{Lp型Vitali判据与无穷远紧性|hyperpage}
有限测度空间上，若 $f_n\to f$ 依测度，则
\[
\|f_n-f\|_p\to0\iff\{|f_n|^p\}\text{ 一致可积}\qquad(1\le p<\infty).
\]
在 $\mathbb R^d$ 上，可改用局部依测度收敛，并同时要求 $L^p$ 有界、$p$ 次密度一致可积及
\[
\lim_{R\to\infty}\sup_n\int_{|x|>R}|f_n(x)|^p\,\mathrm dx=0.
\]
最后一项排除平移导致的质量逃逸；只控制大函数值的尾部不控制空间无穷远。
\section{Sobolev差分、代表与边界}\label{sec:ra-sobolev-tools}
\subsection{差商估计与弱导数判别}\index{差商估计与弱导数判别|hyperpage}
定义 $\delta_h^ju(x)=[u(x+he_j)-u(x)]/h$。全空间 $u\in W^{1,p}$ 满足
\[
\|\delta_h^ju\|_p\le\|\partial_ju\|_p,\qquad
\delta_h^ju\to\partial_ju\text{ 于 }L^p\quad(1\le p<\infty).
\]
反之，$1<p<\infty$ 时，$u\in L^p$ 及小差商的一致 $L^p$ 界蕴含对应弱导数存在。开集内须缩小区域，使差分线段留在域中；$p=1$ 的差商界一般只保证 BV 正则性。
\subsection{沿线绝对连续与截断链式法则}\index{沿线绝对连续与截断链式法则|hyperpage}
$u\in W^{1,p}_{\rm loc}$ 有沿几乎每条坐标平行直线局部绝对连续的代表，且在这些直线上
\[
u(x+te_j)-u(x+se_j)=\int_s^t\partial_ju(x+re_j)\,\mathrm dr.
\]
对 $T_k(t)=\max(-k,\min(t,k))$，$D(T_k(u))=\bm1_{\{|u|<k\}}Du$ a.e.；$D(u^+)=\bm1_{\{u>0\}}Du$。这些法则按弱导数理解，不要求 $u$ 处处经典可微。
\subsection{光滑稠密性与延拓算子}\index{光滑稠密性与延拓算子|hyperpage}
对任意开集及 $1\le p<\infty$，$C^\infty(\Omega)\cap W^{k,p}(\Omega)$ 在 $W^{k,p}(\Omega)$ 中稠密。对有界 Lipschitz 域，存在有界线性延拓
\[
E:W^{1,p}(\Omega)\to W^{1,p}(\mathbb R^d),\quad Eu|_\Omega=u,\quad
\|Eu\|_{W^{1,p}(\mathbb R^d)}\le C\|u\|_{W^{1,p}(\Omega)}.
\]
这里可取 $1\le p\le\infty$。任意函数直接作零延拓一般不保留弱导数；零迹函数的零延拓才有相应结论。
\subsection{迹不等式与零边界刻画}\index{迹不等式与零边界刻画|hyperpage}
有界 Lipschitz 域上，迹算子满足
\[
\|Tu\|_{L^p(\partial\Omega)}\le C\|u\|_{W^{1,p}(\Omega)},\qquad
\ker T=W_0^{1,p}(\Omega)\quad(1\le p<\infty).
\]
对 $1<p<\infty$，精确迹空间为 $W^{1-1/p,p}(\partial\Omega)$，由边界图册定义，并有有界右逆。边界值是函数空间中的迹，不是对任意等价类直接逐点限制。
\section{函数不等式的临界与定量形式}
\subsection{Hardy权重不等式}\index{Hardy权重不等式|hyperpage}
对 $u\in C_c^\infty(\mathbb R^d)$、$1<p<d$，
\[
\int\frac{|u|^p}{|x|^p}\,\mathrm dx
\le\left(\frac p{d-p}\right)^p\int|\nabla u|^p\,\mathrm dx.
\]
常数最优，但非零极值函数不在相应齐次空间中取得该等号。$p=d$ 是不同的临界情形，通常需要对数修正，不能直接让上式中的常数取极限。
\subsection{Gagliardo--Nirenberg插值}\index{Gagliardo-Nirenberg插值|hyperpage}
若 $1\le p<d$、$1\le q\le p^*=dp/(d-p)$、$q\le r\le p^*$，并选 $\theta\in[0,1]$ 满足
\[
\frac1r=\theta\left(\frac1p-\frac1d\right)+\frac{1-\theta}q,
\]
则光滑紧支撑函数满足 $\|u\|_r\le C\|\nabla u\|_p^\theta\|u\|_q^{1-\theta}$。这是 Sobolev 不等式与范数插值的组合；更高阶的完整 Gagliardo--Nirenberg 形式还需检查导数阶数及端点例外。
\subsection{Nash不等式与耗散尺度}\index{Nash不等式与耗散尺度|hyperpage}
对 $u\in H^1(\mathbb R^d)\cap L^1(\mathbb R^d)$，
\[
\|u\|_2^{2+4/d}\le C_d\|\nabla u\|_2^2\|u\|_1^{4/d}.
\]
结合热方程的 $L^1$ 收缩与能量耗散，可解出 $\|u(t)\|_2\le C(\kappa t)^{-d/4}\|u_0\|_1$。该推导连接函数空间估计与半群平滑化，不能把仅有能量有界当作同样的衰减结论。
\subsection{Moser--Trudinger指数可积性}\index{Moser-Trudinger指数可积性|hyperpage}
有界域 $\Omega\subset\mathbb R^d$、$d\ge2$ 上，令 $\alpha_d=d|S^{d-1}|^{1/(d-1)}$。则
\[
\sup_{\substack{u\in W_0^{1,d}(\Omega)\\\|\nabla u\|_d\le1}}
\int_\Omega\exp\bigl(\alpha|u|^{d/(d-1)}\bigr)\,\mathrm dx<\infty
\quad(0\le\alpha\le\alpha_d).
\]
对非空有界开域，$\alpha>\alpha_d$ 时该上确界为无穷。它是临界 $p=d$ 的替代嵌入；零边界条件或相应均值归一化不能删除。
\section{函数族的紧性与集中}\label{sec:ra-compactness}
\subsection{Fréchet--Kolmogorov紧性判据}\index{Fréchet-Kolmogorov紧性判据|hyperpage}
对 $1\le p<\infty$，$\mathcal F\subset L^p(\mathbb R^d)$ 相对紧当且仅当有界，并同时满足
\begin{gather*}
\lim_{R\to\infty}\sup_{f\in\mathcal F}\int_{|x|>R}|f|^p=0,\\
\lim_{h\to0}\sup_{f\in\mathcal F}\|f(\cdot+h)-f\|_p=0.
\end{gather*}
它分别控制无穷远逃逸与小尺度振荡；只要求每个函数自身平移连续不够，必须对整个函数族一致。
\subsection{Rellich嵌入的指数范围}\index{Rellich嵌入的指数范围|hyperpage}
有界 Lipschitz 域上，$1\le p<\infty$ 时，$W^{1,p}(\Omega)$ 紧嵌入 $L^q(\Omega)$ 的常用范围为
\[
1\le q<p^*\ (p<d),\qquad1\le q<\infty\ (p=d),\qquad1\le q\le\infty\ (p>d).
\]
当 $p>d$，还紧嵌入 $C^{0,\beta}(\overline\Omega)$，其中 $0\le\beta<1-d/p$。临界 Hölder 指数及 $q=p^*$ 一般不能保证紧性；域的延拓性质是使用这些全局结论的条件。
\subsection{Aubin--Lions时间紧性}\index{Aubin-Lions时间紧性|hyperpage}
设 Banach 空间 $X_0\Subset X\hookrightarrow X_1$。若 $1\le p<\infty$，函数族满足
\[
\sup_{u\in\mathcal F}\bigl(\|u\|_{L^p(0,T;X_0)}+\|u_t\|_{L^1(0,T;X_1)}\bigr)<\infty,
\]
则在 $L^p(0,T;X)$ 中相对紧。若改为 $L^\infty(0,T;X_0)$ 有界、时间导数在 $L^r(0,T;X_1)$ 有界且 $r>1$，则可得 $C([0,T];X)$ 紧性。时间导数按分布意义理解。
\subsection{集中函数与三种质量行为}\index{集中函数与三种质量行为|hyperpage}
对 $\rho_n\ge0$、$\int\rho_n=M>0$，定义
\[
Q_n(R)=\sup_y\int_{B_R(y)}\rho_n,\qquad
\alpha=\lim_{R\to\infty}\lim_{n\to\infty}Q_n(R),
\]
其中先沿子列在整数半径取得极限。$\alpha=0$ 对应消失；$\alpha=M$ 对应适当平移后的紧性；$0<\alpha<M$ 对应可分离的两团质量。集中紧性允许识别临界问题中的紧性缺失，但不自动给出强收敛或排除二分。
\section{分数阶实域估计}\label{sec:ra-fractional}
\subsection{Gagliardo半范数与差分刻画}\index{Gagliardo半范数与差分刻画|hyperpage}
对 $0<s<1$、$1\le p<\infty$，
\[
[u]_{W^{s,p}(\mathbb R^d)}^p=\iint\frac{|u(x)-u(y)|^p}{|x-y|^{d+sp}}\,\mathrm dx\,\mathrm dy
=\int\frac{\|u(\cdot+h)-u\|_p^p}{|h|^{d+sp}}\,\mathrm dh.
\]
与 $\|u\|_p$ 一起定义 $W^{s,p}$；$p=2$ 时等价于 $H^s$ 范数。对一般 $p\ne2$，不可把 Slobodeckij 空间与 Bessel 势空间无条件认作同一空间。
\subsection{分数嵌入与Hölder范围}\index{分数嵌入与Hölder范围|hyperpage}
全空间上 $sp<d$、$p_s^*=dp/(d-sp)$ 时，
\[
\|u\|_{p_s^*}\le C[u]_{W^{s,p}}\quad(u\in C_c^\infty).
\]
当 $sp>d$，$W^{s,p}$ 具有指数 $s-d/p$ 的 Hölder 代表。对有界 Lipschitz 域，$sp=d$ 时嵌入每个有限 $L^q$；$sp<d$ 时在低于 $p_s^*$ 的指数处紧嵌入。临界端点不能用简单代入得到 $L^\infty$。
\subsection{分数阶迹与边界损失}\index{分数阶迹与边界损失|hyperpage}
有界 Lipschitz 域上，若 $1<p<\infty$、$1/p<s<1$，则
\[
T:W^{s,p}(\Omega)\to W^{s-1/p,p}(\partial\Omega),\qquad
\|Tu\|_{W^{s-1/p,p}(\partial\Omega)}\le C\|u\|_{W^{s,p}(\Omega)}.
\]
边界限制损失 $1/p$ 阶；可用有界右逆恢复边值。$s=1/p$ 的端点需使用更精细的空间条件，不包含在该式中。
\subsection{BBM极限与零阶极限}\index{BBM极限与零阶极限|hyperpage}
对 $u\in W^{1,p}(\mathbb R^d)$、$1\le p<\infty$，
\[
\lim_{s\uparrow1}(1-s)[u]_{W^{s,p}}^p
=K_{d,p}\int|\nabla u|^p,\qquad K_{d,p}=\frac1p\int_{S^{d-1}}|\omega\cdot e|^p\,\mathrm d\sigma.
\]
$e$ 为任意单位向量。若 $u\in W^{s_0,p}(\mathbb R^d)$、$s_0>0$，则 $\lim_{s\downarrow0}s[u]_{W^{s,p}}^p=2|S^{d-1}|\|u\|_p^p/p$。归一化常数随半范数定义改变；有界域的零阶极限不能直接套用全空间公式。
\section{多维BV的结构与边界}\label{sec:ra-bv}
\subsection{导数的三部分分解}\index{导数的三部分分解|hyperpage}
对标量 $u\in BV(\Omega)$，其向量值分布导数为有限测度，并分解为
\[
Du=\nabla u\,\mathrm dx+(u^+-u^-)\nu_u\,
\mathcal H^{d-1}_{\rm norm}\!\restriction J_u+D^cu.
\]
$J_u$ 为近似跳跃集，$u^\pm$ 为两侧近似极限，$\nu_u$ 选择与两侧标号一致的方向，$D^cu$ 为 Cantor 部分。Hausdorff 测度取与平面面积一致的归一化；$SBV$ 恰要求 $D^cu=0$。
\subsection{BV紧性与严格逼近}\index{BV紧性与严格逼近|hyperpage}
有界 Lipschitz 域上，$\sup_n(\|u_n\|_1+|Du_n|(\Omega))<\infty$ 蕴含存在 $L^1$ 收敛子列，且
\[
|Du|(\Omega)\le\liminf_n|Du_n|(\Omega).
\]
对 $u\in BV(\Omega)$，可选光滑 $u_j$ 使 $\|u_j-u\|_1\to0$ 且 $\int_\Omega|\nabla u_j|\to|Du|(\Omega)$，称严格逼近。变差总量收敛不等于导数测度在全变差范数下收敛。
\subsection{BV迹与Gauss--Green公式}\index{BV迹与Gauss-Green公式|hyperpage}
有界 Lipschitz 域上有 $Tu\in L^1(\partial\Omega)$，满足 $\|Tu\|_1\le C(\|u\|_1+|Du|(\Omega))$。对 $\phi\in C^1(\overline\Omega;\mathbb R^d)$，
\[
\int_\Omega u\operatorname{div}\phi\,\mathrm dx+\int_\Omega\phi\cdot\mathrm dDu
=\int_{\partial\Omega}Tu\,\phi\cdot n\,\mathrm d\mathcal H^{d-1}_{\rm norm}.
\]
有限周长集在全空间中的公式为 $\int_E\operatorname{div}\phi=\int_{\partial^*E}\phi\cdot\nu_E\,\mathrm d\mathcal H^{d-1}_{\rm norm}$，$\nu_E$ 为外法向；故 $D\bm1_E=-\nu_E\mathcal H^{d-1}_{\rm norm}\restriction\partial^*E$。
\subsection{BV链式法则与跳跃贡献}\index{BV链式法则与跳跃贡献|hyperpage}
若 $\Phi\in C^1(\mathbb R)$ 且 $\Phi'$ 有界，则局部有
\begin{align*}
D(\Phi(u))&=\Phi'(\widetilde u)(\nabla u\,\mathrm dx+D^cu)\\
&\quad+(\Phi(u^+)-\Phi(u^-))\nu_u\mathcal H^{d-1}_{\rm norm}\restriction J_u.
\end{align*}
$\widetilde u$ 为非跳跃处的近似代表。若要求全局 $\Phi(u)\in L^1$，还需域有限测度或 $\Phi(0)=0$ 等条件。跳跃部分必须使用两侧函数值之差，不能简单用单点 $\Phi'(u)$ 乘以整个 $Du$。
\section{重排与Orlicz可积性}
\subsection{Hardy平均算子}\index{Hardy平均算子|hyperpage}
在 $(0,\infty)$ 上令 $Hf(t)=t^{-1}\int_0^tf(s)\,\mathrm ds$。对 $1<p<\infty$，
\[
\|Hf\|_p\le\frac p{p-1}\|f\|_p.
\]
常数最优。对非负递减重排 $f^*$，$f^{**}=H(f^*)$；该式比较重排平均与原重排，可用于构造 Lorentz 空间的等价范数。$p=1$ 不具有同样的强型界。
\subsection{Pólya--Szegő重排不等式}\index{Pólya-Szegő重排不等式|hyperpage}
非负函数的对称递减重排记 $u^\#$，其超水平集为与 $\{u>t\}$ 等体积的同心球。对 $u\in W^{1,p}(\mathbb R^d)$、$1\le p<\infty$，
\[
\|u^\#\|_p=\|u\|_p,\qquad\|\nabla u^\#\|_p\le\|\nabla u\|_p.
\]
此处各正超水平集有限测度，由 $u\in L^p$ 保证。梯度能量不增连接等周不等式与径向极值问题；一般有边界问题须先确认零延拓或可容许类。
\subsection{Young函数与Luxemburg范数}\index{Young函数与Luxemburg范数|hyperpage}
设 $\Phi$ 为凸递增函数，$\Phi(0)=0$、$\Phi(t)\to\infty$。定义
\[
\|f\|_{L^\Phi}=\inf\left\{\lambda>0:\int\Phi(|f|/\lambda)\,\mathrm d\mu\le1\right\}.
\]
有限范数函数构成 Orlicz 空间。$\Phi(t)=t^p$ 给出 $L^p$；指数增长和 $t\log(1+t)$ 则控制不同的可积性。范数收敛与模积分收敛的相互转换通常还需 $\Delta_2$ 等增长条件。
\subsection{互补Young函数与广义Hölder}\index{互补Young函数与广义Hölder|hyperpage}
令 $\Psi(s)=\sup_{t\ge0}(st-\Phi(t))$，则
\[
st\le\Phi(t)+\Psi(s),\qquad
\int|fg|\,\mathrm d\mu\le2\|f\|_{L^\Phi}\|g\|_{L^\Psi}.
\]
这里使用 Luxemburg 范数，因子 $2$ 与范数约定相关。$\Phi(t)=t^p/p$ 时 $\Psi(s)=s^{p'}/p'$；若讨论整个 Orlicz 对偶，必须另查可分性及 $\Delta_2$ 条件，不能一律只写成互补空间。

\section{精细代表、容量与逐点控制补充}
Sobolev 容量和逐点理论参见\cite{kinnunen}，BV 能量与水平集参见\cite{variation}。
\subsection{Sobolev容量与拟连续代表}\index{Sobolev容量与拟连续代表|hyperpage}
对紧集 $K\subset\mathbb R^d$、$1<p<\infty$，定义
\[
\operatorname{Cap}_p(K)=\inf_{\substack{v\in C_c^\infty\\v\ge1\text{ 于 }K\text{ 邻域}}}
\int_{\mathbb R^d}(|v|^p+|\nabla v|^p).
\]
以通常内外正则方式延拓到一般集合。每个 $W^{1,p}(\mathbb R^d)$ 元有拟连续代表：任意 $\varepsilon>0$ 可删去容量小于 $\varepsilon$ 的开集，使其余处连续；代表在容量零集之外唯一。容量零强于 Lebesgue 零测。
\subsection{Sobolev函数的Hajłasz逐点界}\index{Sobolev函数的Hajłasz逐点界|hyperpage}
若 $u\in W^{1,p}(\mathbb R^d)$、$1<p<\infty$，则在一个零测例外集之外，
\[
|u(x)-u(y)|\le C_d|x-y|
\bigl(M(|\nabla u|)(x)+M(|\nabla u|)(y)\bigr).
\]
因而可取 $g=C_dM(|\nabla u|)\in L^p$ 作为逐点梯度控制，$\|g\|_p\le C_{d,p}\|\nabla u\|_p$。$p=1$ 时极大算子一般只有弱型界，不能保持同样的 $L^1$ 估计。
\subsection{Campanato平均振荡与Hölder连续性}\index{Campanato平均振荡与Hölder连续性|hyperpage}
对 $0<\alpha\le1$、$1\le p<\infty$，令
\[
[u]_{\alpha,p}=\sup_{x\in\mathbb R^d,\,0<r\le1}
r^{-\alpha}\left(\fint_{B(x,r)}|u-u_{B(x,r)}|^p\right)^{1/p}.
\]
若 $u\in L^p_{\rm loc}$ 且此量有限，则有连续代表，满足 $|u(x)-u(y)|\le C[u]_{\alpha,p}|x-y|^\alpha$（$|x-y|\le1/2$）。反向估计亦成立；端点 $\alpha=1$ 对应 Lipschitz，而 $\alpha=0$ 的 BMO 一般不连续。
\subsection{BV余面积公式与水平集周长}\index{BV余面积公式与水平集周长|hyperpage}
实值 $u\in BV(\Omega)$、$\Omega$ 开，有限周长按相对区域解释，则
\[
|Du|(\Omega)=\int_{-\infty}^{\infty}
P(\{u>t\};\Omega)\,\mathrm dt.
\]
几乎每个水平超集有有限周长。该式把分布导数总变差分解为水平集边界能量；当 $u$ 光滑时退化到通常余面积公式，但 BV 情形还允许跳跃。
\clearpage
\end{document}
